% !TeX spellcheck = de_DE

Die Datengrundlage der Analyse basiert auf realen Anwendungsfällen aus dem Nachrichtenbereich des Südwestrundfunks.
Für jeden Anwendungsfall liegt ein Tripel von Texten vor: der ursprüngliche Text in Standardsprache, eine manuell übersetzte Version des Textes in LS+ sowie eine erzeugte Version in LS+, welche aus der bestehenden LS+-Pipeline entstanden ist.
Die automatisch generierten LS+-Texte bilden den Ausgangspunkt für den in dieser Arbeit entwickelten Feedbackloop und werden im Folgenden als Ausgangstexte $T_0$ bezeichnet.

Zentral für die Evaluierung ist die Beziehung zwischen dem durch die Pipeline generierten Text und der im Feedbackloop überarbeiteten Version.

Die pipeline LS+ Version $T_0$ repräsentiert den jetzigen Stand der bestehenden Pipeline und den praktischen Einsatzpunkt der neu entwickelten Komponente.
Die im Rahmen der Evaluation betrachteten Veränderungen der Metriken zwischen Ausgangstext $T_0$ und den Versionen die durch den Feedbackloop generiert wurden $T_1$, $T_2$, ..., $T_n$ sind ausschlaggebend für die Beurteilung, ob der betrachtete Ansatz einen Mehrwert für Texte in LS+ bieten kann.

Beschreiben der Quellen; Was für Texte; Wie viele Texte; Redaktion Barrierefreiheit; Reale Artikel;
Grundlage Beschreiben mit was arbeite ich;


Die Datengrundlage der vorliegenden Untersuchung basiert auf redaktionellen Nachrichtentexten des Südwestrundfunks (SWR), die im Zeitraum von Oktober2025 bis Juni2026 entstanden sind. Für jeden Anwendungsfall liegt ein Tripel von Textvarianten vor: der ursprüngliche Nachrichtentext in Standardsprache, eine automatisch erzeugte Version in Leichter Sprache Plus (LS+) aus der bestehenden LS+-Pipeline sowie eine manuell erstellte LS+-Version. Damit umfasst der verfügbare Korpus insgesamt über \textbf{500 unterschiedliche Nachrichtenthemen}, zu denen jeweils drei Textvarianten vorliegen, sodass sich insgesamt mehr als \textbf{1,500 Einzeltexte} ergeben.

Die Nachrichtenthemen decken einen breiten Ausschnitt des redaktionellen Spektrums ab. Vertreten sind unter anderem Beiträge aus den Bereichen Politik und Gesellschaft (z.,B.\ \enquote{Koalitionsausschuss Bürgergeld}, \enquote{Behinder-ten-gleich-stellungs-gesetz}, \enquote{Wahlprogramm Inklusion}), internationale Politik und Konflikte (z.,B.\ \enquote{Einigung Hamas und Israel}, \enquote{Ukraine}, \enquote{Irankrieg}), Wirtschaft und Arbeit (z.,B.\ \enquote{Chipmangel}, \enquote{Automobilindustrie}, \enquote{Mindestlohn}), Umwelt und Klima (z.,B.\ \enquote{Waldbrände}, \enquote{Weltklimakonferenz}, \enquote{Wasserentnahme}), Gesundheit und Medizin (z.,B.\ \enquote{E-Zigaretten}, \enquote{Darmkrebs}, \enquote{Krankenhaus LU}), Kultur und Freizeit (z.,B.\ \enquote{Comiccon}, \enquote{Schokofestival}, \enquote{Nibelungenfestspiele}) sowie Sport (z.,B.\ \enquote{HandballEM}, \enquote{Paralympics}, \enquote{Fußball-WM}). Darüber hinaus enthält der Korpus zahlreiche Beiträge mit direktem Bezug zu Barrierefreiheit und Inklusion, etwa \enquote{Lebenshilfe}, \enquote{Tag der Leichten Sprache}, \enquote{Inklusionsbetrieb} oder \enquote{Inklusionsbeauftragter}. Diese thematische Breite stellt sicher, dass der entwickelte Ansatz nicht nur an triviale oder künstliche Beispiele, sondern an realistische, redaktionell relevante Szenarien aus dem Nachrichtenalltag des SWR angewendet wird.

Im Mittelpunkt der Evaluierung stehen die automatisch generierten LS+-Texte aus der bestehenden Pipeline. Sie werden im Folgenden als Ausgangstexte 
T
0
T 
0
​
bezeichnet und bilden den praktischen Einsatzpunkt des entwickelten Feedbackloops. Die manuellen LS+-Versionen werden ergänzend für die qualitative Einordnung der Ergebnisse herangezogen, um die vom System erzeugten Überarbeitungen mit redaktionell erstellten Zieltexten vergleichen zu können.

Zur Veranschaulichung der Korpusstruktur dient Abbildung~\ref{fig:korpus-struktur}. Sie zeigt, dass für jeden Anwendungsfall mindestens der Originaltext in Standardsprache und eine automatische LS+-Version vorliegen, die in den Feedbackloop eingespeist wird, während für einen Teil der Fälle zusätzlich eine manuelle LS+-Version als Referenz verfügbar ist.

